\documentclass[10pt,a4paper]{amsart}
\usepackage[margin=19mm]{geometry}
\usepackage{amsmath,amssymb,mathtools}
\usepackage[scr=rsfs,cal=euler]{mathalfa}
\usepackage{xfrac}
\usepackage[unicode,hidelinks,bookmarks=false]{hyperref}
\newcommand{\ie}{i.e.\ }
\newcommand{\cf}{cf.\ }
\newcommand{\resp}{respectively\ }
\newcommand{\Subsection}{\subsection}
\providecommand{\nobreakdash}{\nobreak\hbox{-}}
\setlength{\emergencystretch}{2em}
\multlinegap0pt
\usepackage{amssymb}
\usepackage{xcolor}
\usepackage[shortlabels]{enumitem}
\usepackage{tikz}
\newtheorem{proposition}{Proposition}
\newtheorem{lemma}{Lemma}
\usepackage{cleveref}
\crefname{proposition}{Proposition}{Propositions}
\crefname{lemma}{Lemma}{Lemmas}
\newcommand{\TSsym}{\mathsf{TS}}
\title{\textbf{Token-sliding realizability for complements, Cartesian-products, and grid graph families}}
\author{Duc A. Hoang}
\date{Source: arXiv:2606.03765v1 (CC BY 4.0)}
\begin{document}
\maketitle

\noindent\emph{Source and licence.} \textbf{Token-sliding realizability for complements, Cartesian-products, and grid graph families}. Version arXiv:2606.03765v1 (CC BY 4.0), distributed by arXiv under the Creative Commons Attribution 4.0 International licence, http://creativecommons.org/licenses/by/4.0/, as recorded in the arXiv licence metadata for this exact version and shown on the abstract page https://arxiv.org/abs/2606.03765v1. Full licence notice: \texttt{licenses/arxiv-2606.03765-CC-BY-4.0.txt}.\par\smallskip

\noindent\textit{Editorial scope.} Counted unit: the article's lemma lem:ts2-four-regular-c4-cap (source0.tex lines 1105--1109). The article's complete original proof of it is delivered with it (source0.tex lines 1111--1265, one \texttt{proof} environment of 155 lines). No mathematics is written, repaired or re-derived: the statement and the proof are the article's own lines, in the article's own notation, and no retained line is substituted or re-flowed.  Retained uncounted as the source-local context the proof needs: lines 473--480; lines 1057--1069.  Nothing was omitted from any retained span, so the package carries no omission record.  No retained line contains a citation, so the wrapper carries no bibliography and no bibliography entry.  No retained line contains a figure environment, an image-inclusion command or drawing code, so no image is delivered and the package declares no asset dependency.  Editorial formatting only, in addition: an AMS \texttt{amsart} preamble that restates the article's own declarations and macros used by the retained lines, namely the environments \texttt{proposition}, \texttt{lemma} on separate counters, together with the standard packages those lines need.  The retained environments print in this document as Proposition 1, Lemma 1 in order of appearance, whereas the article prints the same items with the numbering of its own full document.  The complete native source of the article as distributed in the arXiv e-print for this exact version is bundled unchanged as \texttt{sources/201999/source0.tex}.
\begin{proposition}
\label{prop:per-edge-degree}
Let $F$ be a simple graph and let $ab\in E(F)$ be identified with a vertex of $\TSsym_2(\overline{F})$.
Then
\[
 \deg_{\TSsym_2(\overline{F})}(ab)\;=\;\deg_F(a)+\deg_F(b)-2-2t_F(ab).
\]
\end{proposition}

\begin{lemma}
\label{lem:ts2-local-c4-count}
Let $F$ be a graph, let $H:=\TSsym_2(\overline{F})$, and suppose that $H$ is triangle-free.
Fix an edge $ab\in E(F)$, viewed as a vertex of $H$, and put $A:=N_F(a)\setminus N_F[b]$, $B:=N_F(b)\setminus N_F[a]$, and $C:=N_F(a)\cap N_F(b)$.
Let $\mathcal C_H(ab)$ be the set of induced $4$-cycles of $H$ that contain the vertex~$ab$.
Then
\[
|\mathcal C_H(ab)|
=m_F(A,B)
+\sum_{t\in C}\binom{|A\setminus N_F(t)|}{2}
+\sum_{t\in C}\binom{|B\setminus N_F(t)|}{2}.
\]
\end{lemma}

\begin{lemma}
\label{lem:ts2-four-regular-c4-cap}
Let $F$ be a graph and let $H:=\TSsym_2(\overline{F})$.
If $H$ is triangle-free and $4$-regular, then every vertex of $H$ lies on at most four induced $4$-cycles.
\end{lemma}

\begin{proof}
We prove the bound by applying the local four-cycle count with the degree limits from $4$-regularity.
Once the two sides of the edge are denoted by $A$ and $B$, only a few sizes of $A$ and $B$ are possible, and each gives a bound of four.
The cases are organized by the sizes $x=|A|$ and $y=|B|$.
The $(3,1)$ and $(1,3)$ cases are immediate from the side-sum bound.
The one-sided $(4,0)$ and $(0,4)$ cases contain the main counting work, where the possibilities $|C|=6$ and $|C|=5$ are excluded.
The balanced $(2,2)$ case is controlled by a final common-neighbor count.

Fix a vertex $ab\in V(H)$, equivalently an edge $ab\in E(F)$, and use the notation $A,B,C$ from \cref{lem:ts2-local-c4-count}.
Since $H$ is $4$-regular, $|A|+|B|=4$.
For $t\in C$, set $x:=|A|$, $y:=|B|$, $p_A(t):=|A\setminus N_F(t)|$, and $p_B(t):=|B\setminus N_F(t)|$.

The following local degree identity, which is \cref{prop:per-edge-degree} in the present notation, will be used repeatedly.
For every edge $rs\in E(F)$,
\begin{equation}
\label{eq:product-local-edge-degree}
\deg_H(rs)=\deg_F(r)+\deg_F(s)-2-2|N_F(r)\cap N_F(s)|.
\end{equation}
Indeed, the neighbors of $rs$ in $H$ are the edges $rw$ with $w\in N_F(r)\setminus N_F[s]$, together with the edges $sw$ with $w\in N_F(s)\setminus N_F[r]$.

For $t\in C$, subtracting \eqref{eq:product-local-edge-degree} for the two edges $at$ and $bt$ gives
\begin{equation}
\label{eq:product-pa-pb}
p_A(t)-p_B(t)=x-2.
\end{equation}
Here $\deg_F(a)=1+x+|C|$ and $\deg_F(b)=1+y+|C|$.
The common neighbors of $a,t$ are $b$, the vertices of $A\cap N_F(t)$, and the vertices of $C\cap N_F(t)$.
The common neighbors of $b,t$ are $a$, the vertices of $B\cap N_F(t)$, and the vertices of $C\cap N_F(t)$.

For each $u\in A$, the vertex $au$ of $H$ is adjacent to $ab$, to each $at$ with $t\in C\setminus N_F(u)$, and to each $uv$ with $v\in B\cap N_F(u)$.
These are distinct neighbors, and $H$ is $4$-regular; hence $|C\setminus N_F(u)|+|B\cap N_F(u)|\le 3$.
Summing over $u\in A$ yields $\displaystyle m_F(A,B)+\sum_{t\in C}p_A(t)\le 3x$.
Symmetrically,
\begin{equation}
\label{eq:product-b-side-sum}
m_F(A,B)+\sum_{t\in C}p_B(t)\le 3y.
\end{equation}

By \cref{lem:ts2-local-c4-count}, the number of induced $4$-cycles of $H$ through $ab$ is
\begin{equation}
\label{eq:product-q-count}
Q:=m_F(A,B)+\sum_{t\in C}\binom{p_A(t)}2
+\sum_{t\in C}\binom{p_B(t)}2.
\end{equation}
We prove $Q\le 4$ by cases on $(x,y)$.

\begin{itemize}
\item \textbf{Case 1: $(x,y)\in\{(3,1),(1,3)\}$.}
By symmetry, it suffices to treat $(x,y)=(3,1)$.
In this case, \eqref{eq:product-pa-pb} gives $(p_A(t),p_B(t))=(1,0)$ or $(2,1)$.
Hence, the summand $\displaystyle \binom{p_A(t)}2+\binom{p_B(t)}2$ is $p_B(t)$, and \eqref{eq:product-b-side-sum} gives $Q\le 3$.

\item \textbf{Case 2: $(x,y)\in\{(4,0),(0,4)\}$.}
By symmetry, it suffices to treat $(x,y)=(4,0)$.
Then \eqref{eq:product-pa-pb} gives $p_A(t)=2$ for every $t\in C$, so $Q=|C|$.
Let $A=\{u_1,u_2,u_3,u_4\}$, set $c:=|C|$, and set $r_i:=|N_F(u_i)\cap C|$.
Since every vertex of $C$ is adjacent to exactly two vertices of $A$,
\begin{equation}
\label{eq:product-one-sided-rsum}
\sum_{i=1}^4 r_i=2c.
\end{equation}
Applying \eqref{eq:product-local-edge-degree} to $au_i$ gives
\begin{equation}
\label{eq:product-one-sided-deg}
\deg_F(u_i)=7-c+2r_i.
\end{equation}
For $i<j$, let $\lambda_{ij}$ be the number of vertices of $C$ adjacent to both $u_i$ and $u_j$, and let $g_{ij}$ be the number of common neighbors of $u_i,u_j$ outside $\{a,b\}\cup A\cup C$.
The common neighbors of $u_i$ and $u_j$ are $a$, the other two vertices of $A$, the $\lambda_{ij}$ vertices of $C$ adjacent to both, and the $g_{ij}$ common neighbors outside $\{a,b\}\cup A\cup C$.
Thus, \eqref{eq:product-local-edge-degree}, together with \eqref{eq:product-one-sided-deg}, gives
\begin{equation}
\label{eq:product-one-sided-lambda}
\lambda_{ij}+g_{ij}=1-c+r_i+r_j.
\end{equation}
Moreover, each $t\in C$ contributes to exactly one $\lambda_{ij}$, so $\displaystyle \sum_{i<j}\lambda_{ij}=c$.
Summing \eqref{eq:product-one-sided-lambda} over all six pairs and using \eqref{eq:product-one-sided-rsum}, we get $\displaystyle \sum_{i<j}g_{ij}=6-c$.
Thus, $c\le 6$.

We exclude $c=6$ and $c=5$.
If $c=6$, then every $g_{ij}=0$.
From \eqref{eq:product-one-sided-lambda}, we have $\lambda_{ij}=r_i+r_j-5$ for every pair $i<j$.
Using $\displaystyle \sum_i r_i=2c=12$, the identity $\displaystyle r_i=\sum_{j\ne i}\lambda_{ij}$ gives
\[
	r_i=\sum_{j\ne i}(r_i+r_j-5)=3r_i+(12-r_i)-15=2r_i-3.
\]
Hence $r_i=3$ for all $i$, and then $\lambda_{ij}=1$ for every pair~$i<j$.
Equation \eqref{eq:product-one-sided-deg} gives $\deg_F(u_i)=7$, so each $u_i$ has no neighbor outside $\{a\}\cup A\cup C$.

Let $t\in C$, and suppose that $t$ has a neighbor $z$ outside $\{a,b\}\cup A\cup C$.
Since $B=\emptyset$ and $z\notin C$, the outside vertex $z$ is adjacent to neither $a$ nor~$b$.
If $t$ is adjacent to $u_i,u_j\in A$, then the triples $\{b,u_i,z\}$ and $\{b,u_j,z\}$ in $N_F(t)$ force $zu_i,zu_j\in E(F)$; otherwise those triples would give triangles in~$H$.
The adjacency contradicts $g_{ij}=0$.
Hence, no vertex of $C$ has a neighbor outside $\{a,b\}\cup A\cup C$.

Now, each $t\in C$ has neighbors $a,b$, exactly two vertices of $A$, and its remaining neighbors in~$C$.
Applying \eqref{eq:product-local-edge-degree} to $bt$ shows that every $t\in C$ has exactly three neighbors in~$C$.
Also $\alpha(F[C])\le 2$, since an independent triple in $C$ would give a triangle of $H$ on three edges incident with~$b$.
Hence, $\overline{F[C]}$ is triangle-free and $2$-regular on six vertices, so $\overline{F[C]}\cong C_6$.
Thus, $F[C]$ is the triangular prism.

For any edge $tt'\in E(F[C])$, \eqref{eq:product-local-edge-degree} gives $|N_F(t)\cap N_F(t')|=4$.
The vertices $t,t'$ have no outside common neighbor, always share $a,b$, and share at most one vertex of $A$, so they must have a common neighbor in~$C$.
Thus, every edge of $F[C]$ lies in a triangle of~$F[C]$.
Such a triangle containment is impossible in the triangular prism, whose three matching edges lie in no triangle.

If $c=5$, then $\displaystyle \sum_{i<j}g_{ij}=1$.
By \eqref{eq:product-one-sided-deg}, the number of neighbors of $u_i$ outside $\{a,b\}\cup A\cup C$ is $r_i-2$, and hence the total number of incidences from $A$ to outside vertices is $\displaystyle \sum_i(r_i-2)=2c-8=2$.
An outside vertex adjacent to $s$ vertices of $A$ contributes $s$ to this incidence count and $\displaystyle \binom{s}{2}$ to $\displaystyle \sum g_{ij}$.
The total incidence count is $2$, so either one outside vertex is adjacent to two vertices of $A$, or two outside vertices are each adjacent to one vertex of~$A$.
The latter possibility contributes $0$ to $\displaystyle \sum g_{ij}$, contradicting $\displaystyle \sum g_{ij}=1$.
Hence, there is a unique outside vertex $z$ adjacent to exactly two vertices of~$A$.
Say $z$ is adjacent to $u_1,u_2$.
Then $r_1=r_2=3$ and $r_3=r_4=2$,
and $g_{12}=1$ while all other $g_{ij}$ are zero.
Substituting these values in \eqref{eq:product-one-sided-lambda} gives $\lambda_{12}=\lambda_{13}=\lambda_{14}=\lambda_{23}=\lambda_{24}=1$ and $\lambda_{34}=0$.
Let $t_{ij}$ be the unique vertex of $C$ adjacent to $u_i,u_j$, and put $S_1:=\{t_{12},t_{13},t_{14}\}$ and $S_2:=\{t_{12},t_{23},t_{24}\}$.
Let $s_i:=|N_F(z)\cap S_i|$.
No vertex outside $\{a,b\}\cup A\cup C$ is adjacent to $a$ or~$b$, and $u_1$ has no outside neighbor other than~$z$.
Therefore, the common neighbors of $u_1$ and $z$ are precisely $u_2$ together with $N_F(z)\cap S_1$.
Applying \eqref{eq:product-local-edge-degree} to $u_1z$ gives $\deg_F(z)=2s_1$.
Similarly, applying it to $u_2z$ gives $\deg_F(z)=2s_2$.
Hence, $s_1=s_2=:s$.
Since $S_1\cap S_2=\{t_{12}\}$, the vertex $z$ has at least $2s-1$ neighbors in $C$, and therefore $\deg_F(z)\ge 2+(2s-1)=2s+1$,
contradicting $\deg_F(z)=2s$.
Therefore, $c\le 4$, and so $Q=|C|\le 4$.

\item \textbf{Case 3: $(x,y)=(2,2)$.}
Let $A=\{u_1,u_2\}$ and $B=\{v_1,v_2\}$.
By \eqref{eq:product-pa-pb}, each $t\in C$ has $p_A(t)=p_B(t)$.
Let $h$ be the number of vertices $t\in C$ that are anticomplete to $A\cup B$; these are precisely the vertices with $p_A(t)=p_B(t)=2$.
Then \eqref{eq:product-q-count} becomes
\begin{equation}
\label{eq:product-balanced-q}
Q=m_F(A,B)+2h.
\end{equation}
Let $\kappa,\mu,h$ denote respectively the numbers of vertices $t\in C$ with $p_A(t)=p_B(t)=0,1,2$.
Put $r_i:=|N_F(u_i)\cap C|$.
Then
\begin{equation}
\label{eq:product-balanced-rsum}
r_1+r_2=2\kappa+\mu.
\end{equation}
The edge equation for $au_i$ gives
\begin{equation}
\label{eq:product-balanced-deg}
\deg_F(u_i)=5-|C|+2r_i.
\end{equation}
Let $q_A$ be the number of vertices of $B$ adjacent to both $u_1$ and $u_2$, and let $g_A$ be the number of common neighbors of $u_1,u_2$ outside $\{a,b\}\cup A\cup B\cup C$.
The common neighbors of $u_1,u_2$ are $a$, the $q_A$ vertices of $B$, the $\kappa$ vertices of $C$ adjacent to both $u_1,u_2$, and the $g_A$ outside vertices.
Applying \eqref{eq:product-local-edge-degree} to $u_1u_2$, using \eqref{eq:product-balanced-rsum}, \eqref{eq:product-balanced-deg}, and $|C|=\kappa+\mu+h$, gives $q_A+g_A+h=1$.
Thus, $h\le 1$.
If $h=0$, then $Q=m_F(A,B)\le 4$.
If $h=1$, then $q_A=0$, so no vertex of $B$ is adjacent to both vertices of~$A$.
It follows that $m_F(A,B)\le 2$, and \eqref{eq:product-balanced-q} gives $Q\le 2+2=4$.
\end{itemize}
\end{proof}

\end{document}
